\section{Pengertian Game dan Elemen-elemen Game}

\begin{learningoutcomes}
\item Mendefinisikan game secara formal dan akademis
\item Mengidentifikasi elemen-elemen pembentuk game (Mechanics, Dynamics, Aesthetics)
\item Menganalisis hubungan antara mechanics, dynamics, dan aesthetics
\item Menerapkan konsep MDA framework dalam analisis game
\end{learningoutcomes}

\subsection{Definisi Game}

\begin{definisi}[Game]
Game adalah aktivitas terstruktur yang melibatkan satu atau lebih pemain, dengan tujuan yang jelas, aturan yang mengatur interaksi, dan umpan balik yang memberikan informasi tentang kemajuan menuju tujuan \cite{Hunicke2004}.
\end{definisi}

Menurut \cite{Salen2004}, game memiliki karakteristik utama:
\begin{enumerate}
\item \textbf{Tujuan}: Setiap game memiliki tujuan yang harus dicapai pemain
\item \textbf{Aturan}: Aturan yang membatasi bagaimana pemain dapat mencapai tujuan
\item \textbf{Umpan Balik}: Informasi yang diberikan kepada pemain tentang status mereka
\item \textbf{Partisipasi Sukarela}: Pemain memilih untuk bermain dan menerima aturan
\end{enumerate}

\subsection{MDA Framework}

MDA (Mechanics, Dynamics, Aesthetics) merupakan framework analisis yang dikembangkan oleh \cite{Hunicke2004} untuk memahami game.

\subsubsection{Mechanics}

\begin{definisi}[Mechanics]
Mechanics adalah aturan dan sistem yang mengatur bagaimana game berfungsi. Mechanics mendefinisikan apa yang dapat dilakukan pemain dan bagaimana sistem game merespons aksi tersebut \cite{Hunicke2004}.
\end{definisi}

Contoh mechanics:
\begin{itemize}
\item Sistem pergerakan karakter
\item Mekanik pertempuran
\item Sistem inventori
\item Mekanik crafting
\item Sistem leveling dan progression
\end{itemize}

\subsubsection{Dynamics}

\begin{definisi}[Dynamics]
Dynamics adalah perilaku yang muncul dari interaksi mechanics saat game dimainkan. Dynamics adalah hasil dari mechanics yang bekerja bersama \cite{Hunicke2004}.
\end{definisi}

Contoh dynamics:
\begin{itemize}
\item Tension yang muncul dari time limit
\item Kompetisi antar pemain
\item Eksplorasi dan discovery
\item Risk vs reward decisions
\item Emergent gameplay
\end{itemize}

\subsubsection{Aesthetics}

\begin{definisi}[Aesthetics]
Aesthetics adalah pengalaman emosional yang dirasakan pemain saat bermain game. Aesthetics adalah respons emosional terhadap dynamics \cite{Hunicke2004}.
\end{definisi}

Jenis-jenis aesthetics menurut MDA framework:
\begin{enumerate}
\item \textbf{Sensation}: Game sebagai sense-pleasure
\item \textbf{Fantasy}: Game sebagai make-believe
\item \textbf{Narrative}: Game sebagai drama
\item \textbf{Challenge}: Game sebagai obstacle course
\item \textbf{Fellowship}: Game sebagai social framework
\item \textbf{Discovery}: Game sebagai uncharted territory
\item \textbf{Expression}: Game sebagai self-discovery
\item \textbf{Submission}: Game sebagai pastime
\end{enumerate}

\subsection{Hubungan Mechanics-Dynamics-Aesthetics}

\begin{figure}[h]
\centering
\begin{tikzpicture}[node distance=2cm]
    \node [block] (mechanics) {Mechanics\\Aturan Sistem};
    \node [block, below=of mechanics] (dynamics) {Dynamics\\Perilaku yang Muncul};
    \node [block, below=of dynamics] (aesthetics) {Aesthetics\\Pengalaman Emosional};
    
    \draw [line, ->] (mechanics) -- node[right] {Menghasilkan} (dynamics);
    \draw [line, ->] (dynamics) -- node[right] {Menciptakan} (aesthetics);
    \draw [line, ->, bend right=30] (aesthetics) to node[below] {Memengaruhi Desain} (mechanics);
\end{tikzpicture}
\caption{Hubungan Mechanics-Dynamics-Aesthetics}
\label{fig:mda_relationship}
\end{figure}

\subsection{Contoh Analisis Game Menggunakan MDA}

\begin{contoh}
\textbf{Analisis Game Platformer -- Super Mario Bros}

\textbf{Mechanics}: Jump mechanics dengan gravity, collision detection dengan platform, power-up system (mushroom, flower), enemy AI behavior, score system.

\textbf{Dynamics}: Timing jump untuk menghindari musuh, eksplorasi untuk menemukan power-up, risk-reward dalam mengambil coin, tension dari time limit.

\textbf{Aesthetics}: Challenge (mengatasi obstacle), Discovery (menemukan secret area), Sensation (visual dan audio feedback), Fantasy (menjadi karakter hero).
\end{contoh}

\begin{catatan}
MDA framework membantu desainer memahami bahwa perubahan pada mechanics akan mempengaruhi dynamics dan akhirnya aesthetics. Desain game adalah proses iteratif yang memerlukan testing dan refinement.
\end{catatan}
